\section{Bridge: refused primitives and the PQ admission profile}
\label{sec:bridge-pq}

\subsection{Why bridges are different}

The bridge surface is the only Lux surface where the counterparty is not a
Lux subsystem. Every other transcript in the strict-PQ profile is signed by
a key registered in the Z-Chain identity rollup; bridges transact with
external chains, external custodians, and external admission keys.

A strict-PQ chain can choose to refuse classical bridges entirely (the
maximally-paranoid posture: no classical signer is ever admitted), or it can
admit classical bridges under a clearly-labelled, profile-gated path (the
operationally-realistic posture: existing classical infrastructure remains
usable during migration but cannot claim the post-quantum end-to-end label).

LUX\_STRICT\_E2E\_PQ takes the second posture for a specific reason: a
locked profile cannot mandate the world is post-quantum. It can mandate
that \emph{this chain} is post-quantum. The bridge admission profile is
how a strict-PQ chain declares which counterparties it accepts and what
label it claims about itself.

The reference impl is the \texttt{BridgeProfile} in
\texttt{luxfi/bridge@v1.0.7} integrated through
\texttt{luxfi/consensus@v1.23.3}.

\subsection{BridgeProfile struct}

\begin{lstlisting}[language=Go,caption={BridgeProfile fields. Source:
\texttt{luxfi/bridge@v1.0.7}, \texttt{bridge/profile.go}.}]
type BridgeProfile struct {
    ProfileID             BridgeProfileID
    ProfileName           string

    AllowedAdmissionSchemes []SigSchemeID  // schemes admitted for bridge admission

    RefusedPrimitives []RefusedPrimitive   // explicitly-named refused primitives

    AllowsClassicalAdmin   bool  // false on strict-PQ
    PostQuantumEndToEnd    bool  // true iff E2E PQ; gated on AllowsClassicalAdmin=false
    RequireZChainBinding   bool  // every admission carries ZChain anchor
    RequireFinalityProof   bool  // every admission carries Pulsar-M cert

    AdmissionScheme        SigSchemeID    // the admission key's own scheme
    BridgeHashSuiteID      HashSuiteID    // bridge transcript hash family
}
\end{lstlisting}

The strict-PQ instance pins:

\begin{table}[H]
\centering
\caption{Strict-PQ BridgeProfile field values.}
\label{tab:bridge-profile-strict}
\begin{tabular}{ll}
\toprule
Field & Value \\
\midrule
\texttt{ProfileID}              & \texttt{BridgeProfileStrictPQ} \\
\texttt{AllowedAdmissionSchemes} & \{\texttt{ML-DSA-65}, \texttt{ML-DSA-87}\} \\
\texttt{RefusedPrimitives}      & see Table~\ref{tab:bridge-refused} \\
\texttt{AllowsClassicalAdmin}   & \texttt{false} \\
\texttt{PostQuantumEndToEnd}    & \texttt{true} \\
\texttt{RequireZChainBinding}   & \texttt{true} \\
\texttt{RequireFinalityProof}   & \texttt{true} \\
\texttt{AdmissionScheme}        & ML-DSA-65 \\
\texttt{BridgeHashSuiteID}      & SHA3-384 / cSHAKE256 \\
\bottomrule
\end{tabular}
\end{table}

\subsection{Refused primitives}

The refused-primitive list is the affirmative statement of what the strict-PQ
bridge cannot accept:

\begin{table}[H]
\centering
\caption{Refused bridge primitives on the strict-PQ profile.}
\label{tab:bridge-refused}
\begin{tabular}{ll}
\toprule
Primitive & Reason \\
\midrule
ECDSA secp256k1     & Shor breaks the scheme \\
ECDSA secp256r1     & Shor breaks the scheme \\
ed25519             & Shor breaks the scheme \\
ed448               & Shor breaks the scheme \\
BLS12-381 (any agg) & Shor breaks the pairing \\
RSA-2048 / RSA-3072 & Shor breaks integer factorisation \\
BN254 (Groth16 verifier) & Shor breaks pairing; trusted-setup risk \\
KZG (BN254 / BLS12-381)  & Shor breaks pairing; SRS trapdoor risk \\
SHA-1 / SHA-2 (chain mode) & not part of FIPS 202; \texttt{HashSuiteID} mismatch \\
HMAC-SHA-256 (bridge MAC)  & not part of FIPS 202 SHA3 family \\
ECDH X25519 / X448  & Shor breaks Diffie-Hellman \\
\bottomrule
\end{tabular}
\end{table}

The list is exhaustive of the classical primitives that bridges in 2026 use.
Adding a primitive to the list is a profile bump --- the existing strict-PQ
profile stays as-is; a new \texttt{BridgeProfileStrictPQ-v2} variant
includes the newly-refused primitive.

\subsection{AllowsClassicalAdmin and PostQuantumEndToEnd}

The two booleans are not independent. The constraint:

\[
\texttt{PostQuantumEndToEnd} = \texttt{true} \quad \Longrightarrow \quad
\texttt{AllowsClassicalAdmin} = \texttt{false}.
\]

A bridge profile that has \texttt{AllowsClassicalAdmin = true} is, by
definition, \emph{not} post-quantum end-to-end --- a classical admission
signer is the weak link. The profile constructor (\texttt{NewBridgeProfile})
refuses any combination where \texttt{PostQuantumEndToEnd = true} and
\texttt{AllowsClassicalAdmin = true}. The strict-PQ \texttt{Validate}
re-asserts this on every profile load.

This is the formal version of ``you cannot claim PQ E2E if you have a
classical admin key''. The profile cannot \emph{say} the chain is PQ-E2E if
it admits classical primitives.

\subsection{RequireZChainBinding}

Every bridge admission --- every event admitted onto Lux as a
deposit / mint / lock --- must carry a Z-Chain identity anchor: the
admission key is registered in \texttt{ZRegistryRoots.IdentityRoot} at the
admitting epoch. The verifier walks the inclusion proof against the bound
root.

This means a bridge admission key cannot be a wholly-external classical
key: even when the admission is signed off-chain, the signing key's
\texttt{ML-DSA-65} pubkey must be on Z-Chain.

The classical-bridge migration story is therefore: existing bridges
register a new ML-DSA-65 key in the Z-Chain identity rollup, sign all
new admissions under that key, and retire the classical key. The bridge
software runs both keys during the migration window; the strict-PQ profile
admits only the ML-DSA-65-signed events.

\subsection{RequireFinalityProof}

Every admission also carries a Pulsar-M-65 finality certificate
(Section~\ref{sec:pulsar-m}) over the source-chain block height the
admission references. The certificate is verified against the
\texttt{FinalitySchemeID} on the chain admitting the event. For
inter-Lux-chain bridges (C-Chain to X-Chain, etc.), the certificate is the
Q-Chain Pulsar-M cert directly. For external chains, the bridge software
maintains a translator that signs an attestation under ML-DSA-65 over the
external chain's finality artifact; the strict-PQ profile admits the
attestation iff the attestation key is in Z-Chain.

\subsection{RPC label: post\_quantum\_end\_to\_end}

The bridge JSON-RPC exposes a \texttt{bridge\_profile} endpoint that
returns the profile's public claims. The returned object includes a
\texttt{label} field; the strict-PQ profile sets:

\begin{lstlisting}[language=Java,caption={Bridge profile RPC response.}]
{
    "profile_id": "BridgeProfileStrictPQ",
    "profile_hash": "93efd103...323bfcd0",
    "label": "post_quantum_end_to_end",
    "allowed_admission_schemes": ["ML-DSA-65", "ML-DSA-87"],
    "refused_primitives": ["ecdsa-secp256k1", "ed25519", ...],
    "allows_classical_admin": false,
    "post_quantum_end_to_end": true,
    "require_zchain_binding": true,
    "require_finality_proof": true
}
\end{lstlisting}

The \texttt{label} field is a structured public claim. A monitoring service
that scrapes the field is the canonical way to confirm the bridge's
posture; a classical-compat bridge returns
\texttt{label: "classical\_compat"} instead, and the monitor sees the
difference.

\subsection{Permissive profile reciprocal}

The classical-compat bridge profile (\texttt{BridgeProfilePermissive}) is
the reciprocal:

\begin{itemize}[leftmargin=2em,nosep]
  \item \texttt{AllowsClassicalAdmin = true}
  \item \texttt{PostQuantumEndToEnd = false}
  \item \texttt{AllowedAdmissionSchemes} includes both ML-DSA-65 and
    ECDSA-secp256k1
  \item \texttt{RefusedPrimitives} is empty (or short)
\end{itemize}

The permissive profile is used by testnets and by chains in active
migration. The strict-PQ profile loads at the end of the migration window.
Bridges that have not migrated by the cut-over event are not admitted on
the strict-PQ chain; their classical-compat counterparts continue to
operate on testnet / devnet under the permissive profile.
